\section{The target is a jump form}\label{sec:jump}
\begin{proposition}[Jump form]\label{prop:jump}
For $f\in C_c^\infty(I_L)$ and $F=E_Lf$ put
\begin{equation}\label{eq:levy-measure}
 \mu_L=2\gammak(r)\dd r+2\!\!\sum_{m\log p<L}\!\!(\log p)p^{-m/2}\delta_{m\log p},
 \qquad
 \gamma_L=w_0-2\!\!\sum_{m\log p<L}\!\!(\log p)p^{-m/2},
\end{equation}
and
\begin{equation}\label{eq:jump-energy}
 \EL[F]=\int_0^\infty\big[\norm F^2-\re\ip F{U_rF}\big]\dd\mu_L(r).
\end{equation}
Then $\mu_L\geq0$, $\EL\geq0$, and
\begin{equation}\label{eq:jump-decomposition}
 Q_L[f]=\EL[E_Lf]+\gamma_L\norm f^2+P_L[f].
\end{equation}
\end{proposition}
\begin{proof}
For the continuous part of \eqref{eq:levy-measure}, \eqref{eq:jump-energy} is
\eqref{eq:gamma-energy}. For the atomic part, write each prime term of
\eqref{eq:target} as
$-2c\,\re\ip F{U_sF}=-2c\norm F^2+2c[\norm F^2-\re\ip F{U_sF}]$ with
$c=(\log p)p^{-m/2}$ and $s=m\log p$; the first terms collect into
$\gamma_L-w_0$ and the second into the atomic part of \eqref{eq:jump-energy}.
Nonnegativity is $\norm F^2-\re\ip F{U_rF}=\tfrac12\norm{F-U_rF}^2$.
\end{proof}

Equivalently, in symbols, $b(\tau^2)+w_0-2\sum c\cos(\tau s)=\Phi_L(\tau)+\gamma_L$
with $\Phi_L(\tau)=\int(1-\cos\tau r)\dd\mu_L(r)$ a continuous negative definite
function in the sense of Schoenberg. So \eqref{eq:jump-decomposition} is a
L\'evy--Khinchine representation of the whole of \eqref{eq:target} except the
pole term, and it is elementary; what is worth saying is what it means.

\paragraph{One object, not two channels.} The archimedean density and the prime
atoms are the continuous and atomic parts of the single nonnegative measure
$\mu_L$. A source for \eqref{eq:target} is, up to the constant and the two
poles, the Dirichlet form of a L\'evy jump process whose jumps are the primes,
with mass $2(\log p)p^{-m/2}$ at jump length $m\log p$. This matters because the
companion manuscript proves that no coherent source has the archimedean energy
and the summed prime references as its two channel norms: the interference bound
fails at explicit test functions and by a growing factor. A presentation in
which the two are one object is what that exclusion asks for, and
\eqref{eq:jump-decomposition} supplies it.

\paragraph{The constant.} $\gamma_L$ is large and negative, $-6.3524$ at $L=1$
and $-18.3878$ at $L=3$, with $\gamma_L=w_0-4e^{L/2}(1+o(1))$ by the prime number
theorem. It is a different bookkeeping from the constant $c_L$ of
\eqref{eq:sigma}, which uses the mirror contact $2dr/(1+r)$ in place of the full
$2dr^n$.

\paragraph{A ground-state reformulation.} Taking the least Rayleigh quotient of
\eqref{eq:jump-decomposition} against $\norm f^2$ gives, exactly,
\begin{equation}\label{eq:ground-state}
 \lambda_{\min}\big(\EL+P_L;\norm\cdot^2\big)
 =-\gamma_L+\lambda_{\min}\big(Q_L;\norm\cdot^2\big),
\end{equation}
so that Weil positivity on $I_L$ says precisely that the ground-state energy of
the jump form together with the pole term is at least
$2\sum_{m\log p<L}(\log p)p^{-m/2}-w_0$. In a $48$-cell basis at $L=1$ the two
sides of \eqref{eq:ground-state} are $6.354438$ and $6.352442$, the difference
being $\lambda_{\min}(Q_L)$. This is a re-presentation and not a theorem, but it
places the criterion where the spectral theory of nonlocal Dirichlet forms on an
interval is the relevant apparatus.
